%
% scene_id_namespacing.tex
%
% A Z specification of DES-086's connection-scoped store-key composition: the
% Hub never uses a client-submitted scene/frame id as a store key directly: it
% composes ConnectionScopedId.compose(connection_id, local_id) at the single
% write choke point (SceneInstaller.install, mirrored at SceneOperations.update
% and SceneClearer.clear). This closes the bug vox hit running in two Claude
% Code sessions at once: both sessions pushed scene_id="music-player" and the
% second silently evicted the first, because the Hub stored every scene under
% the literal string a client submitted, with no ownership check on the first
% write to a given id.
%
% The operator required this model as an override of both the design author's
% and the security reviewer's "not required" recommendation (Decision 4,
% docs/architecture/scene-id-namespacing-design.md), given the project's
% history of missed interleavings on DES-037/038. This model proves the
% design's own "no new interleaving" argument by exhaustive search rather than
% by taking it on the design's word: composition is a pure function of the
% caller's already-resolved ConnectionId, so no shared mutable state is
% contended, and the model-check confirms no operation --- including a
% connection's own socket re-establishing (Reconnect) --- can misattribute a
% key once written.
%
% Composing the store key as the literal pair (connection_id, local_id) is a
% faithful abstraction of a separator-delimited string join: the reserved
% separator's whole job is to make concatenation behave like pairing ---
% unambiguously decomposable --- which is exactly what a Z pair already is by
% construction, given a local_id that cannot itself carry the separator (the
% \texttt{taint} predicate below models that guard). The fidelity control,
% docs/scene_id_namespacing_buggy.tex, drops the connection component from the
% key entirely --- "composition omitted", the pre-DES-086 defect --- and the
% identical goal check that is unreachable here becomes reachable there.
%
% Models the design in docs/architecture/scene-id-namespacing-design.md as
% realised by src/punt_lux/domain/hub/connection_scoped_id.py
% (ConnectionScopedId.compose, the __post_init__ separator/blank guard) and
% src/punt_lux/operations/scene_installer.py (the one choke point every scene
% install passes through).
%
% Type-check:  fuzz -t docs/scene_id_namespacing.tex
% Model-check: probcli docs/scene_id_namespacing.tex -model_check \
%                -p DEFAULT_SETSIZE 3 -p MAX_INITIALISATIONS 200
%   Invariant 1 (collision-impossibility) and Invariant 3 (taint rejected
%   before composition) are reachability goals, not state invariants (so the
%   fidelity variant, which reuses the identical goal shape, can reach the bad
%   state it is built to reproduce):
%     probcli docs/scene_id_namespacing.tex -model_check \
%       -goal "#(c1,c2,l).(c1:CONNECTION & c2:CONNECTION & l:LOCAL_ID & \
%              c1/=c2 & (c1|->l):dom(store) & store(c1|->l)/=c1)" \
%       -p DEFAULT_SETSIZE 3 -p MAX_INITIALISATIONS 200
%     probcli docs/scene_id_namespacing.tex -model_check \
%       -goal "#(c,l).(c:CONNECTION & l:LOCAL_ID & (c|->l):dom(store) & \
%              taint(l)=tainted)" \
%       -p DEFAULT_SETSIZE 3 -p MAX_INITIALISATIONS 200
%   Both goals: never satisfied over the whole reachable state space.
%   Invariant 2 (idempotency) is a corollary of compose being a pure function
%   of (connection_id, local_id) and is argued in Section~\ref{sec:idempotency}
%   rather than searched for --- there is no state-space content in "the same
%   inputs produce the same output" to exhaustively check.
%
% Formatting follows the fuzz house rules: schema lines under 80 columns,
% predicates broken at \land/\lor/\implies.
%

\documentclass[a4paper,10pt,fleqn]{article}
\usepackage[margin=1in]{geometry}
\usepackage{fuzz}
\usepackage[colorlinks=true,linkcolor=blue,citecolor=blue,urlcolor=blue]{hyperref}

\begin{document}

\title{lux Connection-Scoped Store-Key Composition: a Z Specification}
\author{Formal model of DES-086's collision-impossibility invariant}
\date{August 2026}
\maketitle

% ============================================================
\section{Introduction}
% ============================================================

Every scene or frame a client shows is installed under a store key. Before
DES-086, that key was the literal string the client submitted --- its
\texttt{scene\_id} --- with no ownership check on the first write to a given
id. Two different connections that happened to choose the identical string did
not error; the second silently took over as owner, evicting the first's
elements wholesale. This is the bug vox hit running in two Claude Code sessions
at once: both pushed \texttt{scene\_id="music-player"}, and the second overwrote
the first on the shared display.

DES-086's fix composes the store key from the writing connection's own
\texttt{ConnectionId} and the caller's local label,
\texttt{ConnectionScopedId.compose(connection\_id, local\_id)}, at the one
choke point every scene install passes through. A malformed \texttt{local\_id}
--- blank, or carrying the reserved unit separator that joins the two halves
--- is rejected before composition ever runs, so a client cannot construct a
\texttt{local\_id} that would spell another connection's composed key.

Three properties must hold:

\begin{description}
  \item[I1 --- collision-impossibility.] Two distinct connections never end up
    owning the same composed key. Equivalently: the composed key recorded for
    a connection's own write is always attributed to that connection, never to
    another one, across every reachable sequence of writes and reconnects.
  \item[I2 --- idempotency.] The same connection composing the same
    \texttt{local\_id} twice always produces the same key, so a re-show
    updates the one scene it already owns rather than creating a second.
  \item[I3 --- reject before compose.] A \texttt{local\_id} that carries the
    reserved separator is refused before it is ever composed into a key ---
    the unrepresentable input is a boundary error, not a value the model ever
    has to reason about downstream.
\end{description}

The model abstracts the store key as the pair
$(\mathit{connection}, \mathit{local\_id})$ rather than as a literal joined
string. This is a faithful representation, not a simplification that begs the
question: the separator's entire purpose is to make string concatenation
behave like pairing --- unambiguously decomposable back into its two halves ---
which is exactly what a Z pair already is by construction, \emph{given} that
\texttt{local\_id} cannot itself carry the separator. That precondition is
what I3 states and what the \texttt{taint} predicate below encodes. What the
model-check verifies is not the trivial fact that distinct pairs are distinct
--- it is that no operation in the system, across any reachable interleaving,
ever writes a connection's own key with someone else's value, and that no
operation ever composes a key from a \texttt{local\_id} the separator guard
should have refused.

% ============================================================
\section{Basic Types}
% ============================================================

\begin{zed}
  [CONNECTION, LOCAL\_ID]
\end{zed}

\texttt{CONNECTION} is the set of connections the Hub can distinguish ---
\texttt{ConnectionId} values, one per live MCP session or applet connection.
\texttt{LOCAL\_ID} is the set of caller-chosen labels a \texttt{show} call
might submit, before any validation or composition. Two connections already
exhibit the collision question (do connection $c_1$'s and $c_2$'s composed
keys ever collapse); ProB's \texttt{DEFAULT\_SETSIZE} bounds the carrier and a
third connection confirms the invariant holds pairwise, not merely for one
privileged pair.

% ============================================================
\section{Free Types}
% ============================================================

\begin{zed}
  TAINT ::= clean | tainted
\end{zed}

Whether a raw \texttt{local\_id} carries the reserved unit separator ---
\texttt{ID\_SEPARATOR} in the implementation, a control character no honest
caller-chosen id ever contains. \texttt{clean}: composable. \texttt{tainted}:
refused before composition, mirroring \texttt{ConnectionScopedId.\_\_post\_\allowbreak
init\_\_}'s guard.

% ============================================================
\section{Given Functions}
% ============================================================

\begin{axdef}
  taint : LOCAL\_ID \fun TAINT
\end{axdef}

A total, given classification of every \texttt{local\_id} in the carrier ---
the model does not construct actual character sequences (as the fuzz/ProB
toolchain has no native string algebra convenient for this), so which ids are
tainted is fixed axiomatically rather than derived. \texttt{DEFAULT\_SETSIZE}
guarantees ProB enumerates both a \texttt{clean} and a \texttt{tainted}
instantiation of \texttt{LOCAL\_ID} across its search, so the guard is
genuinely exercised, not vacuously true.

% ============================================================
\section{State}
% ============================================================

\begin{schema}{Store}
  store : (CONNECTION \cross LOCAL\_ID) \pfun CONNECTION
\end{schema}

\texttt{store} is the Hub's authoritative map from a composed key --- the pair
--- to the connection that installed it. I1 and I3 are deliberately
\emph{absent} from this schema's predicate: were either a state invariant,
ProB would treat any operation that broke it as \emph{disabled} rather than as
a violation, and the fidelity check in the buggy sibling file could no longer
reach the bad state it exists to reproduce. Both are instead stated as
reachability goals in Section~\ref{sec:verification} and checked against the
full reachable state space.

% ============================================================
\section{Initialization}
% ============================================================

\begin{schema}{Init}
  Store'
\where
  store' = \emptyset
\end{schema}

No scene has been shown; the store holds no keys.

% ============================================================
\section{Operations}
% ============================================================

\texttt{Show} is a caller installing a scene: it composes its own connection
with its own chosen \texttt{local\_id} and writes the pair, attributed to
itself. This is the whole of \texttt{SceneInstaller.install}'s composition
step, modeled: no operation in this system ever writes a key attributed to a
connection other than the one performing the write, because \texttt{Show}
takes only its own $c?$ as both the key's first component and the value ---
there is no operation that could spell "install under a key naming a
connection other than the caller's own." Guarded on \texttt{clean}, so a
tainted \texttt{local\_id} is refused before it ever reaches the store (I3).

\begin{schema}{Show}
  \Delta Store \\
  c? : CONNECTION \\
  l? : LOCAL\_ID
\where
  taint~l? = clean \\
  store' = store \oplus \{ (c?, l?) \mapsto c? \}
\end{schema}

\texttt{Reconnect} models a connection's socket re-establishing --- the DES-068
scenario named in the mission's own framing, ``composition against a
ConnectionId mid-transition.'' A reconnect does not change which
\texttt{ConnectionId} the session composes with (the identifier is stable for
the connection's lifetime, derived once from the MCP transport's session key
or the declared-identity hash); it leaves every already-composed key exactly
as it was. Modeled as the identity transition on \texttt{store}, so the
model-check exercises interleavings of \texttt{Show} and \texttt{Reconnect}
freely and confirms a reconnect never disturbs an existing attribution.

\begin{schema}{Reconnect}
  \Delta Store \\
  c? : CONNECTION
\where
  store' = store
\end{schema}

% ============================================================
\section{Verification}
\label{sec:verification}
% ============================================================

I1 is checked as a reachability goal: the negating state
\[
  B_1 \defs \exists c_1, c_2 : CONNECTION; l : LOCAL\_ID @
    c_1 \neq c_2 \land (c_1, l) \in \dom store \land store~(c_1, l) \neq c_1
\]
--- some entry keyed by $c_1$'s own composed pair ended up attributed to a
different connection --- must be unreachable. In ProB:

\begin{verbatim}
probcli docs/scene_id_namespacing.tex -model_check \
  -goal "#(c1,c2,l).(c1:CONNECTION & c2:CONNECTION & l:LOCAL_ID & \
         c1/=c2 & (c1|->l):dom(store) & store(c1|->l)/=c1)" \
  -p DEFAULT_SETSIZE 3 -p MAX_INITIALISATIONS 200
\end{verbatim}

Reported never satisfied over the whole reachable state space: no sequence of
\texttt{Show} and \texttt{Reconnect} calls, in any interleaving, ever
attributes $c_1$'s own key to $c_2$. This holds because the pair $(c_1, l)$
and the pair $(c_2, l)$ are distinct domain elements the instant $c_1 \neq
c_2$, so \texttt{Show}'s override at $(c_2, l)$ never touches $(c_1, l)$ ---
the exact property that makes the vox collision unrepresentable rather than
merely checked.

I3 is checked the identical way: the negating state
\[
  B_3 \defs \exists c : CONNECTION; l : LOCAL\_ID @
    (c, l) \in \dom store \land taint~l = tainted
\]
--- a tainted \texttt{local\_id} nonetheless reached the store --- must be
unreachable:

\begin{verbatim}
probcli docs/scene_id_namespacing.tex -model_check \
  -goal "#(c,l).(c:CONNECTION & l:LOCAL_ID & (c|->l):dom(store) & \
         taint(l)=tainted)" \
  -p DEFAULT_SETSIZE 3 -p MAX_INITIALISATIONS 200
\end{verbatim}

Reported never satisfied: \texttt{Show}'s precondition $taint~l? = clean$
disables the operation entirely for a tainted $l?$, so no tainted
\texttt{local\_id} ever composes into a key. The bare
\texttt{-model\_check} additionally confirms every operation is covered and
the structural state predicate is preserved across the reachable space.

\paragraph{Deadlock.} There is nothing here to deadlock: \texttt{Show} and
\texttt{Reconnect} touch one piece of state under no lock at all in this
model, because the design's own claim --- confirmed, not merely repeated,
by this exercise --- is that composition is a pure, synchronous computation
using data already resolved before the request begins, contending with
nothing. Whenever a connection and a clean local id exist, \texttt{Show} is
enabled, so the reachable state space has no stuck state; \texttt{Reconnect}
is unconditionally enabled. \texttt{-cbc\_deadlock} finds no state with every
operation disabled.

% ============================================================
\section{Idempotency (I2)}
\label{sec:idempotency}
% ============================================================

I2 --- the same connection composing the same \texttt{local\_id} twice
produces the same key --- is a corollary of \texttt{Show}'s own definition
rather than a claim with reachability content to search for. $\mathit{Show}$
writes at the pair $(c?, l?)$; firing it twice with identical $c?$ and $l?$
computes the identical pair both times, so the second call re-writes the
identical domain entry with the identical value $c?$ --- a semantic no-op from
the store's perspective, which is exactly the "a re-show updates the one scene
it already owns" behavior \texttt{HubDisplay.replace\_scene} already gives a
repeated \texttt{show()} on an unchanged \texttt{scene\_id}. There is no
interleaving or state-space question here: a mathematical function of two
fixed arguments returns one fixed result on every call, by definition of
being a function, and \texttt{ConnectionScopedId.compose} is exactly that ---
a total function of $(connection\_id, local\_id)$ with no hidden state. Unlike
I1 and I3, forcing this through \texttt{-model\_check} would search a state
space with nothing left to find.

\end{document}
